%% ============================================================
%%  Chapter 2 — System Architecture
%%  02_architecture.tex
%% ============================================================
\chapter{System Architecture}
\label{chap:architecture}

\chapterintro{%
  This chapter dissects the internal architecture of the Clinic Management
  API layer by layer. It traces the complete journey of an HTTP request from
  network ingress to JSON response, documents every middleware class in the
  pipeline with its actual PHP source code, explains the routing strategy,
  and maps the project directory to the architectural responsibility of each
  directory. Engineers extending the system should read this chapter in full
  before touching any controller, service, or middleware class.
}

%% ─────────────────────────────────────────────────────────────────────────
\section{Architectural Philosophy}
\label{sec:arch-philosophy}
%% ─────────────────────────────────────────────────────────────────────────

The application is structured around a strict \textbf{layered architecture}
with one-directional dependency flow. Each layer knows about the layer
directly below it and nothing else. This constraint prevents the codebase
from collapsing into a ball of mud as it grows, and it is what makes the
73-test PHPUnit suite possible: each layer can be tested in isolation by
mocking the layer beneath it.

The dependency direction is:

\begin{center}
  \large
  \textbf{HTTP Request}
  $\;\longrightarrow\;$
  \textbf{Middleware}
  $\;\longrightarrow\;$
  \textbf{Controller}
  $\;\longrightarrow\;$
  \textbf{Service}
  $\;\longrightarrow\;$
  \textbf{Repository}
  $\;\longrightarrow\;$
  \textbf{Eloquent / MySQL}
\end{center}

\vspace{6pt}

Controllers depend on Services and Form Requests. Services depend on
Repositories. Repositories depend on Eloquent models. At no point does
a Repository reference a Service, or a Service reference a Controller.
Violations of this constraint are architectural defects and should be
treated as bugs.

The application also follows the \textbf{Explicit Over Implicit} principle.
Every action that modifies data passes through a Form Request (explicit
validation), a Service (explicit business logic), and a Repository
(explicit query). There are no raw \texttt{DB::} calls in controllers,
no business logic in controllers, and no Eloquent queries in controllers.

%% ─────────────────────────────────────────────────────────────────────────
\section{The Laravel 12 Request Lifecycle}
\label{sec:request-lifecycle}
%% ─────────────────────────────────────────────────────────────────────────

Understanding how Laravel 12 processes an HTTP request is essential context
before examining the application-specific layers. The following description
traces the path of a request to
\route{POST /api/v1/assistant/appointments} from the moment it hits the
server to the moment the JSON response is sent.

\subsection{Stage 1 --- Web Server and PHP-FPM}

The web server (Nginx in the recommended production configuration) receives
the raw TCP connection. For any request that does not match a static file,
Nginx passes the request to PHP-FPM via the FastCGI protocol. PHP-FPM
spawns or reuses a PHP worker process, which begins executing
\texttt{public/index.php} --- the application's single entry point.

\subsection{Stage 2 --- Bootstrap and Application Kernel}

\texttt{public/index.php} performs three actions:

\begin{enumerate}
  \item Loads the Composer autoloader (\texttt{vendor/autoload.php}),
    making every class in the application and its dependencies
    discoverable by PHP's class-loading mechanism.
  \item Calls \texttt{require \_\_DIR\_\_.'/../bootstrap/app.php'}, which
    instantiates the Laravel \texttt{Application} container, registers all
    service providers listed in \texttt{bootstrap/providers.php}, and
    configures the routing, middleware, and exception-handling pipelines
    via the fluent builder returned by \texttt{Application::configure()}.
  \item Calls \texttt{\$app->handleRequest(Request::capture())} on the
    configured application instance, which captures the PHP superglobals
    (\texttt{\$\_GET}, \texttt{\$\_POST}, \texttt{\$\_SERVER}) into a
    Symfony \texttt{Request} object and passes it into the HTTP kernel.
\end{enumerate}

\subsection{Stage 3 --- The Global Middleware Stack}

Before the request reaches any route, the HTTP kernel passes it through
the global middleware stack. For API requests, the relevant global
middleware layers are:

\begin{itemize}
  \item \textbf{TrimStrings} and \textbf{ConvertEmptyStringsToNull} ---
    sanitise all incoming string values.
  \item \textbf{HandleCors} --- evaluates and appends CORS headers if
    cross-origin requests are configured.
  \item \textbf{SubstituteBindings} --- resolves route model binding,
    turning the \texttt{\{id\}} segment in a URL into a concrete Eloquent
    model instance via \texttt{findOrFail()}.
\end{itemize}

\subsection{Stage 4 --- Route Matching}

The router inspects the HTTP method and URL path against the registered
route definitions in \texttt{routes/api.php}. All API routes are prefixed
with \texttt{api/} by the framework and with \texttt{v1/} by the
application's own route group (see Section~\ref{sec:routing}).

If no route matches, Laravel throws a
\texttt{Symfony\textbackslash{}Component\textbackslash{}Http\textbackslash{}Kernel\textbackslash{}Exception\textbackslash{}NotFoundHttpException},
which the exception handler renders as a \stnotfound{} JSON response.

\subsection{Stage 5 --- Route-Specific Middleware}

Once a route is matched, the HTTP kernel runs the middleware stack
attached to that route or route group. For the assistant appointment
routes, this stack is (in execution order):

\begin{enumerate}
  \item \texttt{auth:api} --- Laravel Passport validates the bearer token.
    If the token is absent, expired, or revoked, a
    \texttt{AuthenticationException} is thrown, which the exception
    handler renders as \stunauth{}.
  \item \texttt{role:assistant} --- The custom
    \phpclass{RoleMiddleware} verifies that the authenticated user's
    role enum value equals \texttt{"assistant"}. If not, a \stforbid{}
    JSON response is returned immediately and the request does not
    proceed further.
  \item \texttt{clinic.scope} --- The custom
    \phpclass{ClinicScopeMiddleware} injects the authenticated user's
    \texttt{clinic\_id} into the request's data bag, preventing
    clinic-scope forgery.
\end{enumerate}

\subsection{Stage 6 --- Form Request Resolution and Validation}

Before the controller method is called, Laravel's dependency injection
container resolves the type-hinted \texttt{FormRequest} subclass from
the controller's method signature. The Form Request's \texttt{authorize()}
method is called first (returning \texttt{true} for all requests that
have already passed middleware), then its \texttt{rules()} method
provides the validation rules. If validation fails, Laravel throws a
\texttt{ValidationException} automatically, which the exception handler
renders as a \stunproc{} response with field-level error messages.

\subsection{Stage 7 --- Controller, Service, and Repository Execution}

The controller method receives the resolved Form Request and any
route-model-bound Eloquent instances. It delegates immediately to a
Service class. The Service applies business logic and calls one or more
Repository methods. The Repository executes the Eloquent query against
MySQL.

\subsection{Stage 8 --- Resource Transformation and Response}

The controller wraps the model or collection returned by the service in
a Laravel \texttt{JsonResource} (or \texttt{JsonResource::collection()})
and passes it to the \texttt{ApiResponse::success()} trait method. This
method calls \texttt{\$data->toArray(request())} on the resource and
wraps the result in the standard JSON envelope before returning an
\texttt{Illuminate\textbackslash{}Http\textbackslash{}JsonResponse}.

\subsection{Stage 9 --- Exception Handling}

If any stage from 1 through 8 throws an exception, PHP unwinds the call
stack back to the HTTP kernel, which passes the exception to the
application's exception handler registered in \texttt{bootstrap/app.php}.
The handler matches the exception type against a priority-ordered list of
\texttt{\$exceptions->render()} callbacks and returns the appropriate JSON
response.

The complete request lifecycle is summarised in Table~\ref{tab:lifecycle}.

\begin{table}[h!]
  \centering
  \caption{Summary of the HTTP request lifecycle for an API endpoint.}
  \label{tab:lifecycle}
  \renewcommand{\arraystretch}{1.5}
  \begin{tabularx}{\linewidth}{@{}C{0.5cm} L{4cm} X@{}}
    \toprule
    \rowcolor{PrimaryBlue}
    \tblhdr{\#} & \tblhdr{Stage} & \tblhdr{Responsibility} \\
    \midrule
    1 & Web Server / PHP-FPM   & TCP accept, FastCGI dispatch \\
    \rowalt
    2 & Bootstrap / Kernel     & Autoload, DI container, provider registration \\
    3 & Global Middleware      & CORS, string sanitisation, binding substitution \\
    \rowalt
    4 & Route Matching         & Method + path $\rightarrow$ controller action \\
    5 & Route Middleware       & \texttt{auth:api}, \texttt{role:X}, \texttt{clinic.scope} \\
    \rowalt
    6 & Form Request           & Validation, automatic 422 on failure \\
    7 & Controller / Service / Repository & Business logic, database I/O \\
    \rowalt
    8 & Resource + ApiResponse & Model $\rightarrow$ JSON envelope \\
    9 & Exception Handler      & Catch any throwable, return structured error JSON \\
    \bottomrule
  \end{tabularx}
\end{table}

%% ─────────────────────────────────────────────────────────────────────────
\section{Project Directory Structure}
\label{sec:directory}
%% ─────────────────────────────────────────────────────────────────────────

The application follows the standard Laravel directory convention with
additional sub-directories within \texttt{app/} to enforce layer
separation. The annotated tree below shows every directory that contains
application code.

\begin{lstlisting}[
  style     = bash,
  caption   = {Annotated project directory tree.},
  label     = {lst:directory},
  numbers   = none,
]
CMS/
|-- app/
|   |-- Enums/                  # PHP 8.1 backed enums (type-safe constants)
|   |   |-- AppointmentStatus.php
|   |   |-- Gender.php
|   |   `-- UserRole.php
|   |
|   |-- Exceptions/             # Custom domain exception classes
|   |   |-- AccountDeactivatedException.php
|   |   |-- AppointmentConflictException.php
|   |   |-- ClinicScopeViolationException.php
|   |   |-- InvalidAppointmentStateException.php
|   |   `-- SelfDemotionException.php
|   |
|   |-- Http/
|   |   |-- Controllers/        # Thin HTTP adapters; no business logic
|   |   |   |-- Auth/
|   |   |   |   `-- AuthController.php
|   |   |   |-- Assistant/
|   |   |   |   |-- AppointmentController.php
|   |   |   |   `-- PatientController.php
|   |   |   |-- Doctor/
|   |   |   |   |-- AppointmentController.php
|   |   |   |   |-- PatientController.php
|   |   |   |   |-- PrescriptionController.php
|   |   |   |   `-- ScheduleController.php
|   |   |   `-- SuperAdmin/
|   |   |       |-- ClinicController.php
|   |   |       |-- DashboardController.php
|   |   |       `-- UserController.php
|   |   |
|   |   |-- Middleware/         # Cross-cutting request interceptors
|   |   |   |-- ClinicScopeMiddleware.php
|   |   |   `-- RoleMiddleware.php
|   |   |
|   |   |-- Requests/           # FormRequest validation classes
|   |   |   |-- Appointment/
|   |   |   |-- Auth/
|   |   |   |-- Clinic/
|   |   |   |-- Patient/
|   |   |   |-- Prescription/
|   |   |   |-- Schedule/
|   |   |   `-- User/
|   |   |
|   |   `-- Resources/          # JSON resource transformers
|   |       |-- AppointmentResource.php
|   |       |-- ClinicResource.php
|   |       |-- PatientResource.php
|   |       |-- PrescriptionItemResource.php
|   |       |-- PrescriptionResource.php
|   |       |-- ScheduleResource.php
|   |       `-- UserResource.php
|   |
|   |-- Jobs/                   # Async queue jobs
|   |   `-- SendAppointmentConfirmation.php
|   |
|   |-- Models/                 # Eloquent models (HasUuids on all)
|   |   |-- Appointment.php
|   |   |-- Clinic.php
|   |   |-- DoctorSchedule.php
|   |   |-- Patient.php
|   |   |-- Prescription.php
|   |   |-- PrescriptionItem.php
|   |   `-- User.php
|   |
|   |-- Providers/
|   |   `-- AppServiceProvider.php
|   |
|   |-- Repositories/           # All Eloquent query logic lives here
|   |   |-- AppointmentRepository.php
|   |   |-- ClinicRepository.php
|   |   |-- PatientRepository.php
|   |   |-- PrescriptionRepository.php
|   |   |-- ScheduleRepository.php
|   |   `-- UserRepository.php
|   |
|   |-- Services/               # Domain business logic
|   |   |-- AppointmentService.php
|   |   |-- AuthService.php
|   |   |-- ClinicService.php
|   |   `-- PrescriptionService.php
|   |
|   `-- Traits/
|       `-- ApiResponse.php     # Shared success/error response helpers
|
|-- bootstrap/
|   `-- app.php                 # Application config, middleware, exceptions
|
|-- database/
|   |-- factories/              # Model factories for seeding and tests
|   |-- migrations/             # Ordered schema migration files
|   `-- seeders/
|
|-- routes/
|   `-- api.php                 # All API route definitions
|
`-- tests/
    |-- ApiTestCase.php         # Abstract base class for all feature tests
    `-- Feature/
        |-- Auth/
        |-- Assistant/
        |-- Doctor/
        `-- SuperAdmin/
\end{lstlisting}

Table~\ref{tab:dir-responsibilities} maps each directory to its
architectural role and documents the constraints applied to code in
that directory.

\begin{longtable}{@{}L{3.8cm} L{2.8cm} X@{}}
  \toprule
  \rowcolor{PrimaryBlue}
  \tblhdr{Directory} & \tblhdr{Layer} & \tblhdr{Constraints} \\
  \midrule
  \endfirsthead
  \toprule
  \rowcolor{PrimaryBlue}
  \tblhdr{Directory (cont.)} & \tblhdr{Layer} & \tblhdr{Constraints} \\
  \midrule
  \endhead
  \bottomrule
  \endfoot

  \texttt{app/Enums}            & Type definitions
    & No logic. Backed enums only. No imports outside the PHP stdlib. \\
  \rowalt
  \texttt{app/Exceptions}       & Exception definitions
    & Each class \texttt{extends Exception}, defines a default message,
      and implements a \texttt{render()} method returning a JSON response. \\
  \texttt{app/Http/Controllers} & HTTP adapter
    & No Eloquent queries. No \texttt{DB::} calls. No business rules.
      Delegates immediately to Service or Repository. \\
  \rowalt
  \texttt{app/Http/Middleware}  & Request pipeline
    & May read from \texttt{Request} and \texttt{User}. May return early
      with a response or call \texttt{\$next(\$request)}. \\
  \texttt{app/Http/Requests}    & Validation
    & Defines \texttt{rules()}, \texttt{authorize()}, and optionally
      \texttt{messages()}. No direct DB access. \\
  \rowalt
  \texttt{app/Http/Resources}   & Serialisation
    & Transforms a model into a JSON-safe array. No queries. \\
  \texttt{app/Jobs}             & Async work
    & Implements \texttt{ShouldQueue}. Receives model or value objects.
      Does not depend on the HTTP request object. \\
  \rowalt
  \texttt{app/Models}           & ORM wrappers
    & Declares \texttt{\$fillable}, \texttt{\$casts}, relationships,
      and the \texttt{HasUuids} trait. No business logic. \\
  \texttt{app/Repositories}     & Data access
    & All Eloquent queries live here. Called by Services only.
      Returns model instances or collections. \\
  \rowalt
  \texttt{app/Services}         & Business logic
    & Calls Repositories. Throws domain Exceptions. Dispatches Jobs.
      No direct HTTP concern. \\
  \texttt{app/Traits/ApiResponse} & Response helpers
    & Provides \texttt{success()} and \texttt{error()} methods.
      Used only by Controllers. \\
\end{longtable}

%% ─────────────────────────────────────────────────────────────────────────
\section{The Routing Layer}
\label{sec:routing}
%% ─────────────────────────────────────────────────────────────────────────

All API routes are defined in \texttt{routes/api.php}. Laravel
automatically prefixes every route in this file with \texttt{/api},
so the application adds the \texttt{v1} version prefix itself.

\subsection{Route Group Strategy}

Routes are organised into four top-level groups, each protecting its
routes with a combination of the \texttt{auth:api}, \texttt{role}, and
\texttt{clinic.scope} middleware. The grouping strategy serves two
purposes: it eliminates middleware repetition and it makes the
permission model visually obvious to any engineer reading the file.

\begin{lstlisting}[
  style   = php,
  caption = {routes/api.php --- complete route file structure.},
  label   = {lst:routes},
]
<?php

use Illuminate\Support\Facades\Route;

// ── Public Routes (no auth required) ────────────────────────
Route::prefix('v1')->group(function () {

    Route::prefix('auth')->controller(AuthController::class)->group(function () {
        Route::post('register', 'register');
        Route::post('login',    'login');
        Route::middleware('auth:api')->group(function () {
            Route::post('logout', 'logout');
            Route::get('me',      'me');
        });
    });

    // ── Super Admin Routes ───────────────────────────────────
    Route::prefix('admin')
        ->middleware(['auth:api', 'role:super_admin'])
        ->group(function () {

        Route::apiResource('clinics', ClinicController::class);
        Route::apiResource('users',   UserController::class);
        Route::patch('users/{user}/toggle', [UserController::class, 'toggle']);
        Route::patch('users/{user}/role',   [UserController::class, 'updateRole']);
        Route::get('dashboard',             [DashboardController::class, 'index']);
    });

    // ── Doctor Routes ────────────────────────────────────────
    Route::prefix('doctor')
        ->middleware(['auth:api', 'role:doctor', 'clinic.scope'])
        ->group(function () {

        Route::apiResource('schedules',    ScheduleController::class)
            ->except(['show']);
        Route::apiResource('appointments', AppointmentController::class)
            ->only(['index', 'show']);
        Route::patch('appointments/{id}/status',
            [AppointmentController::class, 'updateStatus']);
        Route::apiResource('prescriptions', PrescriptionController::class)
            ->except(['destroy']);
        Route::get('patients',       [PatientController::class, 'index']);
        Route::get('patients/{patient}', [PatientController::class, 'show']);
    });

    // ── Assistant Routes ─────────────────────────────────────
    Route::prefix('assistant')
        ->middleware(['auth:api', 'role:assistant', 'clinic.scope'])
        ->group(function () {

        Route::get('available-slots', [AppointmentController::class, 'availableSlots']);
        Route::apiResource('appointments', AppointmentController::class)
            ->except(['show'])->only(['index', 'store', 'update', 'destroy']);
        Route::get('appointments/{appointment}',
            [AppointmentController::class, 'show']);
        Route::apiResource('patients', PatientController::class)
            ->except(['destroy']);
    });
});
\end{lstlisting}

\subsection{Route Model Binding}

Laravel's implicit route model binding is active for all routes that
include a route segment whose name matches an Eloquent model's
lowercase class name (e.g.\ \texttt{\{clinic\}}, \texttt{\{user\}},
\texttt{\{patient\}}). The framework calls \texttt{findOrFail()} on
the model using the value of the path segment as the primary key lookup.
Since all primary keys are UUIDs, the path segment is a UUID string.

If no record with that UUID exists, Laravel throws a
\texttt{ModelNotFoundException}, which the exception handler registered
in \texttt{bootstrap/app.php} intercepts and renders as a \stnotfound{}
JSON response with the message \texttt{"Resource not found."}.

\begin{notebox}[Route Naming]
  The \texttt{apiResource()} helper registers up to five named routes
  with the convention \texttt{resource.index}, \texttt{resource.store},
  \texttt{resource.show}, \texttt{resource.update}, and
  \texttt{resource.destroy}. The \texttt{->only()} and \texttt{->except()}
  fluent methods limit which of the five routes are generated, ensuring
  that no unintended HTTP method is exposed for a given resource.
\end{notebox}

\subsection{Controller Namespacing}

The route file uses fully-qualified class names imported at the top via
\texttt{use} statements (omitted from Listing~\ref{lst:routes} for
brevity). Each controller is namespaced to its role group:
\phpclass{App\textbackslash{}Http\textbackslash{}Controllers\textbackslash{}SuperAdmin},
\phpclass{App\textbackslash{}Http\textbackslash{}Controllers\textbackslash{}Doctor}, and
\phpclass{App\textbackslash{}Http\textbackslash{}Controllers\textbackslash{}Assistant}.
This means both \phpclass{Doctor\textbackslash{}PatientController} and
\phpclass{Assistant\textbackslash{}PatientController} can coexist without
naming conflicts. Each resolves to a distinct class with distinct
capability boundaries.

%% ─────────────────────────────────────────────────────────────────────────
\section{The Middleware Pipeline}
\label{sec:middleware}
%% ─────────────────────────────────────────────────────────────────────────

The application uses three middleware classes directly. Two are provided
by the Laravel/Passport ecosystem (\texttt{auth:api} and
\texttt{SubstituteBindings}); two are custom application classes.

\begin{table}[h!]
  \centering
  \caption{Complete middleware registry for the Clinic Management API.}
  \label{tab:middleware}
  \renewcommand{\arraystretch}{1.5}
  \begin{tabularx}{\linewidth}{@{}L{3.2cm} L{2.8cm} X@{}}
    \toprule
    \rowcolor{PrimaryBlue}
    \tblhdr{Alias / Class} & \tblhdr{Source} & \tblhdr{Responsibility} \\
    \midrule
    \texttt{auth:api}
      & Laravel Passport
      & Validates the OAuth2 bearer token. Populates
        \texttt{request()->user()} with the authenticated
        \texttt{User} model. \\
    \rowalt
    \texttt{role:\{role\}}
      & Custom (\phpclass{RoleMiddleware})
      & Compares the authenticated user's \texttt{role} enum value
        against the required role string. Returns \stforbid{} if they
        do not match. \\
    \texttt{clinic.scope}
      & Custom (\phpclass{ClinicScopeMiddleware})
      & For non-Super-Admin users, merges the authenticated user's
        \texttt{clinic\_id} into the request data bag, overwriting any
        client-supplied value. \\
    \rowalt
    \texttt{SubstituteBindings}
      & Laravel Framework
      & Resolves route model binding: replaces UUID path segments with
        Eloquent model instances. \\
    \texttt{TrimStrings}
      & Laravel Framework
      & Trims leading and trailing whitespace from all string input fields. \\
    \rowalt
    \texttt{ConvertEmptyStrings ToNull}
      & Laravel Framework
      & Converts empty string input fields to \texttt{null}. \\
    \bottomrule
  \end{tabularx}
\end{table}

\subsection{Middleware Registration in bootstrap/app.php}

Laravel 12 uses a fluent builder in \texttt{bootstrap/app.php} to
register middleware aliases. The \texttt{->withMiddleware()} callback
receives a \texttt{Middleware} builder instance on which the
\texttt{alias()} method maps short string names to fully-qualified
class names.

\begin{lstlisting}[
  style   = php,
  caption = {bootstrap/app.php --- middleware alias registration.},
  label   = {lst:bootstrap-middleware},
]
->withMiddleware(function (Middleware $middleware): void {
    $middleware->alias([
        'clinic.scope' => App\Http\Middleware\ClinicScopeMiddleware::class,
        'role'         => App\Http\Middleware\RoleMiddleware::class,
    ]);
})
\end{lstlisting}

Once registered, these aliases can be used in route group definitions
as short string tokens (e.g.\ \texttt{'role:doctor'}), where the
colon-separated suffix is passed as a runtime parameter to the
middleware's \texttt{handle()} method.

\subsection{RoleMiddleware --- Full Source and Analysis}

The \phpclass{RoleMiddleware} is the application's role-based access
control (RBAC) enforcement point. It is called on every request to the
\texttt{admin}, \texttt{doctor}, and \texttt{assistant} route groups.

\begin{lstlisting}[
  style   = php,
  caption = {app/Http/Middleware/RoleMiddleware.php --- full source.},
  label   = {lst:role-middleware},
]
<?php

namespace App\Http\Middleware;

use Closure;
use Illuminate\Http\Request;
use Symfony\Component\HttpFoundation\Response;

class RoleMiddleware
{
    /**
     * Handle an incoming request.
     *
     * @param  Closure(Request): (Response)  $next
     */
    public function handle(Request $request, Closure $next, string $role): Response
    {
        $user = $request->user();

        // Guard: unauthenticated users are handled upstream by auth:api,
        // but a null check here prevents a fatal error if middleware is
        // applied without auth:api in the stack.
        if (! $user || $user->role->value !== $role) {
            return response()->json([
                'success' => false,
                'message' => 'Forbidden: You do not have the required permissions.'
            ], 403);
        }

        return $next($request);
    }
}
\end{lstlisting}

\textbf{Design analysis.} Several deliberate decisions are embedded in
this twelve-line class.

\begin{enumerate}
  \item \textbf{The \texttt{\$role} parameter is a plain string.}
    The route definition \texttt{'role:super\_admin'} causes Laravel to
    call \texttt{handle(\$request, \$next, 'super\_admin')}. The
    middleware does not know which routes it is attached to; it only
    knows the role it has been asked to enforce. This makes the class
    reusable across all three role groups without modification.

  \item \textbf{Comparison is against \texttt{role->value}, not
    \texttt{role} itself.}
    The \texttt{User} model's \texttt{role} attribute is cast to the
    \phpclass{UserRole} enum by Eloquent. Directly comparing a
    \texttt{UserRole} enum instance to the string \texttt{"doctor"}
    would always return \texttt{false}. The \texttt{->value} accessor
    extracts the underlying string representation from the enum
    (\texttt{"super\_admin"}, \texttt{"doctor"}, or
    \texttt{"assistant"}), making the string comparison correct.

  \item \textbf{The null check on \texttt{\$user} is defensive.}
    In the normal route stack, \texttt{auth:api} runs before
    \texttt{role}, so \texttt{\$request->user()} is always non-null
    when \phpclass{RoleMiddleware} executes. The null check exists to
    prevent a fatal \texttt{Attempt to read property ``role'' on null}
    error in the event that the middleware is accidentally applied
    without the \texttt{auth:api} guard.

  \item \textbf{Early return with \texttt{response()->json()}.}
    The middleware does not throw an exception. It returns a response
    directly. This is slightly different from the custom domain
    exceptions used elsewhere in the application and means that this
    specific 403 response bypasses the central exception-handler
    rendering pipeline. The response format (\texttt{success} +
    \texttt{message}) is manually kept consistent with the
    \texttt{ApiResponse} trait.
\end{enumerate}

\subsection{ClinicScopeMiddleware --- Full Source and Analysis}

The \phpclass{ClinicScopeMiddleware} is the first of the two
data-isolation enforcement layers (the second being the explicit
ownership checks in individual controllers).

\begin{lstlisting}[
  style   = php,
  caption = {app/Http/Middleware/ClinicScopeMiddleware.php --- full source.},
  label   = {lst:scope-middleware},
]
<?php

namespace App\Http\Middleware;

use App\Enums\UserRole;
use Closure;
use Illuminate\Http\Request;
use Symfony\Component\HttpFoundation\Response;

class ClinicScopeMiddleware
{
    /**
     * Handle an incoming request.
     *
     * @param  Closure(Request): (Response)  $next
     */
    public function handle(Request $request, Closure $next): Response
    {
        $user = $request->user();

        // Super Admins operate globally; they bypass all clinic restrictions.
        if ($user->role === UserRole::SUPER_ADMIN) {
            return $next($request);
        }

        // For Doctors and Assistants, forcefully inject their clinic_id
        // into the request data bag. This overwrites any clinic_id value
        // the client may have included in their JSON body, preventing
        // cross-clinic data injection attacks.
        $request->merge([
            'clinic_id' => $user->clinic_id
        ]);

        return $next($request);
    }
}
\end{lstlisting}

\textbf{Design analysis.}

\begin{enumerate}
  \item \textbf{The Super Admin bypass is explicit.}
    The condition \texttt{\$user->role === UserRole::SUPER\_ADMIN}
    compares an enum instance to an enum case --- a type-safe comparison
    that cannot accidentally match a string. Super Admins are routed
    through the \texttt{admin} prefix group, which does not include
    \texttt{clinic.scope} in its middleware stack, so in practice this
    branch is never hit. It exists as a safety net.

  \item \textbf{\texttt{merge()} overwrites, it does not append.}
    \texttt{Request::merge()} adds keys to the request's input data
    bag (accessible via \texttt{request()->all()} and
    \texttt{\$request->clinic\_id}). If the client already sent a
    \texttt{clinic\_id} key in their JSON payload, \texttt{merge()}
    overwrites it with the authenticated user's real value. The client
    cannot supply a different clinic's ID to gain access to its data.

  \item \textbf{Downstream access via \texttt{\$request->clinic\_id}.}
    After this middleware runs, any downstream controller, Form Request,
    or service can read \texttt{\$request->clinic\_id} as if it were
    a first-class request field. This is used by repository methods such
    as \texttt{PatientRepository::allForClinic(\$request->clinic\_id)}
    to automatically filter their queries to the correct clinic.

  \item \textbf{No response is returned early.}
    Unlike \phpclass{RoleMiddleware}, the scope middleware always calls
    \texttt{\$next(\$request)} and never short-circuits the pipeline.
    Clinic-scope violations are discovered later, at the controller
    level, when the explicit ownership check finds that a fetched
    model's \texttt{clinic\_id} does not match the injected value.
\end{enumerate}

%% ─────────────────────────────────────────────────────────────────────────
\section{The Controller Layer}
\label{sec:controller-layer}
%% ─────────────────────────────────────────────────────────────────────────

Controllers in this application are deliberately thin. They perform
exactly three tasks: receive the validated request, delegate to a
service or repository, and return a JSON response. They contain no
business logic, no Eloquent queries, and no conditional branching on
business rules.

\subsection{The ApiResponse Trait}

Every controller uses the \phpclass{ApiResponse} trait, which provides
two methods: \texttt{success()} and \texttt{error()}. These methods
wrap the return value in the standard JSON envelope described in
Section~\ref{sec:response-envelope}.

\begin{lstlisting}[
  style   = php,
  caption = {app/Traits/ApiResponse.php --- the response envelope trait.},
  label   = {lst:api-response},
]
<?php

namespace App\Traits;

use Illuminate\Http\JsonResponse;
use Illuminate\Http\Resources\Json\JsonResource;
use Illuminate\Http\Resources\Json\AnonymousResourceCollection;

trait ApiResponse
{
    protected function success(
        mixed  $data    = null,
        string $message = 'Success',
        int    $status  = 200,
    ): JsonResponse {
        $resolved = match (true) {
            $data instanceof JsonResource,
            $data instanceof AnonymousResourceCollection => $data->toArray(request()),
            is_null($data)                               => null,
            default                                      => $data,
        };

        return response()->json([
            'success' => true,
            'message' => $message,
            'data'    => $resolved,
        ], $status);
    }

    protected function error(
        string $message = 'An error occurred',
        int    $status  = 400,
        array  $errors  = [],
    ): JsonResponse {
        $body = ['success' => false, 'message' => $message];

        if (! empty($errors)) {
            $body['errors'] = $errors;
        }

        return response()->json($body, $status);
    }
}
\end{lstlisting}

\begin{warningbox}[Paginated Collections and the \texttt{meta} Key]
  When \texttt{success()} receives an \texttt{AnonymousResourceCollection}
  backed by a \texttt{LengthAwarePaginator}, the call to
  \texttt{toArray(request())} resolves the paginated data into a flat
  PHP array. The pagination metadata (\texttt{current\_page},
  \texttt{total}, \texttt{last\_page}) is \emph{not} included in the
  response envelope. Client code must not rely on a \texttt{meta} key.
  Pagination is controlled via \texttt{?page=N} query parameters;
  the API returns the items for the requested page in the \texttt{data}
  array.
\end{warningbox}

\subsection{A Representative Controller: AssistantAppointmentController}

The following listing shows the complete
\phpclass{Assistant\textbackslash{}AppointmentController}. It is the
most complex controller in the application and illustrates all the
patterns used throughout the codebase.

\begin{lstlisting}[
  style   = php,
  caption = {app/Http/Controllers/Assistant/AppointmentController.php.},
  label   = {lst:assistant-appt-ctrl},
]
<?php

namespace App\Http\Controllers\Assistant;

use App\Exceptions\ClinicScopeViolationException;
use App\Http\Controllers\Controller;
use App\Http\Requests\Appointment\AvailableSlotsRequest;
use App\Http\Requests\Appointment\StoreAppointmentRequest;
use App\Http\Requests\Appointment\UpdateAppointmentRequest;
use App\Http\Resources\AppointmentResource;
use App\Models\Appointment;
use App\Repositories\AppointmentRepository;
use App\Repositories\PatientRepository;
use App\Services\AppointmentService;
use App\Traits\ApiResponse;
use Illuminate\Http\JsonResponse;
use Illuminate\Http\Request;

class AppointmentController extends Controller
{
    use ApiResponse;

    public function __construct(
        private AppointmentRepository $repo,
        private AppointmentService    $service,
        private PatientRepository     $patientRepo,
    ) {}

    public function index(Request $request): JsonResponse
    {
        $filters      = array_merge(
            ['date' => today()->format('Y-m-d')],
            $request->only(['status', 'date'])
        );
        $clinicId     = $request->user()->clinic_id;
        $appointments = $this->repo->allForClinic($clinicId, $filters);

        return $this->success(
            data: AppointmentResource::collection($appointments)
        );
    }

    public function store(StoreAppointmentRequest $request): JsonResponse
    {
        $patient = $this->patientRepo->getPatientById($request->patient_id);

        if ($patient->clinic_id !== $request->user()->clinic_id) {
            throw new ClinicScopeViolationException(
                'Patient does not belong to this clinic.'
            );
        }

        $data = array_merge($request->validated(), [
            'clinic_id' => $request->user()->clinic_id,
            'booked_by' => $request->user()->id,
        ]);

        $appointment = $this->service->book($data);

        return $this->success(
            data:    new AppointmentResource($appointment),
            message: 'Appointment booked successfully',
            status:  201
        );
    }

    public function show(Appointment $appointment): JsonResponse
    {
        return $this->success(
            data: new AppointmentResource($this->repo->find($appointment->id))
        );
    }

    public function update(UpdateAppointmentRequest $request,
                           Appointment $appointment): JsonResponse
    {
        $updated = $this->service->reschedule($appointment, $request->validated());

        return $this->success(
            data:    new AppointmentResource($updated),
            message: 'Appointment rescheduled successfully'
        );
    }

    public function destroy(Appointment $appointment): JsonResponse
    {
        $this->service->cancel($appointment);

        return $this->success(message: 'Appointment cancelled successfully');
    }

    public function availableSlots(AvailableSlotsRequest $request): JsonResponse
    {
        $slots = $this->service->getAvailableSlots(
            $request->doctor_id,
            $request->date
        );

        return $this->success(data: ['slots' => $slots]);
    }
}
\end{lstlisting}

Key observations about this controller:

\begin{itemize}
  \item \textbf{Constructor injection.} The three dependencies
    (\phpclass{AppointmentRepository}, \phpclass{AppointmentService},
    \phpclass{PatientRepository}) are injected by Laravel's DI container.
    The controller never calls \texttt{new} on any of them.
  \item \textbf{No Eloquent queries.} Every data access goes through a
    repository method. The controller does not call any static Eloquent
    methods such as \texttt{Appointment::where()} directly.
  \item \textbf{No business rules in conditionals.} The only conditional
    in the controller is the clinic-scope ownership check, which throws a
    domain exception rather than returning an inline response.
  \item \textbf{Named arguments throughout.} The \texttt{success()} calls
    use PHP 8.0 named arguments (\texttt{data:}, \texttt{message:},
    \texttt{status:}) to make the intent of each parameter explicit.
\end{itemize}

%% ─────────────────────────────────────────────────────────────────────────
\section{The Service Layer}
\label{sec:service-layer}
%% ─────────────────────────────────────────────────────────────────────────

Service classes are the home of all domain business logic. They are the
only classes permitted to throw the custom domain exceptions defined in
\texttt{app/Exceptions}. They are the only classes permitted to dispatch
queue jobs. They are the only place where multi-step operations
(e.g.\ validate slot, compute end time, clear cache, create record,
dispatch job) are orchestrated into a single transactional unit of work.

\subsection{AppointmentService --- Slot Booking with Cache Lock}

The \phpclass{AppointmentService} is the most complex service in the
application. Its \texttt{book()} method implements a two-phase
concurrency-safe booking algorithm.

\begin{lstlisting}[
  style   = php,
  caption = {app/Services/AppointmentService.php --- book() method.},
  label   = {lst:appt-service-book},
]
public function book(array $data): Appointment
{
    $doctorId = $data['doctor_id'];
    $date     = $data['appointment_date'];
    $start    = $data['start_time'];
    $lockKey  = "booking:{$doctorId}:{$date}:{$start}";

    /** @var Appointment $appointment */
    $appointment = Cache::lock($lockKey, 10)->block(5, function ()
        use ($data, $doctorId, $date, $start)
    {
        // Phase 1: Re-check availability inside the lock.
        // Between the client's slot enquiry and now, another request
        // may have booked the same slot. The lock ensures only one
        // booking attempt proceeds at a time for this slot.
        $availableSlots = $this->getAvailableSlots($doctorId, $date);

        if (! in_array($start, $availableSlots)) {
            throw new AppointmentConflictException();
        }

        // Phase 2: Compute the end time from the schedule's slot_duration.
        $dayOfWeek = Carbon::parse($date)->dayOfWeek;
        $schedule  = $this->scheduleRepo->findForDayAndDoctor($doctorId, $dayOfWeek);

        $data['end_time'] = Carbon::parse($start)
            ->addMinutes($schedule->slot_duration)
            ->format('H:i:s');

        // Invalidate the slot cache so subsequent slot queries reflect
        // this booking immediately.
        Cache::forget($this->slotsCacheKey($doctorId, $date));

        return $this->appointmentRepo->create($data);
    });

    // Dispatch the confirmation job outside the lock to avoid holding
    // the lock while the job is being serialised and queued.
    SendAppointmentConfirmation::dispatch($appointment);

    return $appointment;
}
\end{lstlisting}

\textbf{Why \texttt{Cache::lock()::block()?}}
\texttt{Cache::lock(\$key, \$ttl)} acquires an atomic Redis (or database)
lock. The \texttt{->block(\$timeout, \$callback)} call blocks for up to
\texttt{\$timeout} seconds waiting to acquire the lock, then executes the
callback exclusively. This prevents the following race condition:

\begin{enumerate}
  \item Client A queries available slots. Slot \texttt{08:00} is available.
  \item Client B queries available slots. Slot \texttt{08:00} is available.
  \item Client A submits a booking for \texttt{08:00}.
  \item Client B submits a booking for \texttt{08:00}.
  \item Without a lock, both bookings succeed because both read an empty
    slot before either write completes.
\end{enumerate}

With the lock, step 4 blocks while step 3's lock is held. When step 3
releases the lock and Client B's callback runs, it re-checks slot
availability and finds \texttt{08:00} is now booked. It throws
\phpclass{AppointmentConflictException}, returning \stconflict{}.

%% ─────────────────────────────────────────────────────────────────────────
\section{The Repository Layer}
\label{sec:repository-layer}
%% ─────────────────────────────────────────────────────────────────────────

Repository classes encapsulate every database query in the application.
They depend on Eloquent models but return typed results (model instances
or Eloquent collections). Callers (services) work with the returned
PHP objects and never construct their own queries.

The application uses concrete repositories without interfaces. This is
a deliberate pragmatic decision: the codebase does not have a
requirement to swap database backends, and adding interfaces would double
the number of files without adding testability value, since Laravel's
test infrastructure (\texttt{RefreshDatabase} with a real MySQL instance)
tests the actual repository behaviour end-to-end.

\subsection{AppointmentRepository --- Booked Slot Normalisation}
\label{sec:time-normalisation}

The repository contains a subtle but critical normalisation in the
\texttt{getBookedTimeSlots()} method. MySQL's \texttt{TIME} column type
returns values over the wire as \texttt{HH:MM:SS}
(e.g.\ \texttt{"08:00:00"}). The slot generator in
\phpclass{AppointmentService} formats candidate slots as \texttt{HH:MM}
(e.g.\ \texttt{"08:00"}) using Carbon's \texttt{format('H:i')}. Without
normalisation, the \texttt{in\_array(\$slotTime, \$bookedTimes)} check
would never find a match, making every slot appear available even when
booked.

\begin{lstlisting}[
  style   = php,
  caption = {AppointmentRepository --- MySQL TIME normalisation.},
  label   = {lst:repo-normalise},
]
public function getBookedTimeSlots(string $doctorId, string $date): array
{
    return Appointment::where('doctor_id', $doctorId)
        ->whereDate('appointment_date', $date)
        ->where('status', '!=', 'cancelled')
        ->pluck('start_time')
        // Normalise HH:MM:SS (MySQL) -> HH:MM (slot generator format)
        ->map(fn(string $t) => substr($t, 0, 5))
        ->toArray();
}
\end{lstlisting}

The \texttt{->map(fn(string \$t) => substr(\$t, 0, 5))} call truncates
the seconds component unconditionally. This is safe because MySQL
\texttt{TIME} values in this domain never carry a non-zero seconds
component (all slots are whole-minute boundaries), and \texttt{substr}
of a string shorter than 5 characters returns the full string without
error.

%% ─────────────────────────────────────────────────────────────────────────
\section{The Resource Layer}
\label{sec:resource-layer}
%% ─────────────────────────────────────────────────────────────────────────

Laravel API Resources transform Eloquent model instances into JSON-safe
arrays with explicit field control. Every model in the application has a
corresponding resource class. Resources serve three purposes:

\begin{enumerate}
  \item \textbf{Field whitelisting.} A resource explicitly lists every
    field it exposes. Fields not listed (e.g.\ \texttt{password},
    internal audit flags) are never included in the response, even if
    the underlying model carries them.
  \item \textbf{Relationship expansion.} Eager-loaded relationships are
    conditionally included using the \texttt{whenLoaded()} helper. If the
    relationship was not eager-loaded, the key is omitted from the
    response entirely rather than triggering a lazy-load query.
  \item \textbf{Type coercion.} Enum values, Carbon date instances, and
    integer booleans are cast to their appropriate JSON representation
    within the resource, ensuring consistent output types regardless of
    how the model was constructed.
\end{enumerate}

%% ─────────────────────────────────────────────────────────────────────────
\section{The Job Queue System}
\label{sec:jobs}
%% ─────────────────────────────────────────────────────────────────────────

The application ships with one queue job:
\phpclass{SendAppointmentConfirmation}. This job is dispatched by
\phpclass{AppointmentService::book()} immediately after every successful
appointment creation.

The job is designed to be \emph{idempotent}: if it fails and is retried
by the queue worker, re-executing its logic produces the same observable
outcome (a single confirmation notification). This property is
maintained by ensuring the job does not modify any database state; it
only sends an outbound notification.

In the test environment, the queue driver is set to \texttt{sync},
which executes the job immediately in-process rather than dispatching
it to a real queue. This prevents tests from depending on a running
queue worker and avoids background process leakage between test cases.

\begin{lstlisting}[
  style   = php,
  caption = {app/Jobs/SendAppointmentConfirmation.php.},
  label   = {lst:appt-job},
]
<?php

namespace App\Jobs;

use App\Models\Appointment;
use Illuminate\Contracts\Queue\ShouldQueue;
use Illuminate\Foundation\Queue\Queueable;

class SendAppointmentConfirmation implements ShouldQueue
{
    use Queueable;

    public function __construct(
        public readonly Appointment $appointment
    ) {}

    public function handle(): void
    {
        // Notification dispatch logic:
        // $this->appointment->patient->notify(...)
        // Implementation depends on the configured notification channel.
    }
}
\end{lstlisting}

%% ─────────────────────────────────────────────────────────────────────────
\section{Bootstrap and Exception Handler Configuration}
\label{sec:bootstrap}
%% ─────────────────────────────────────────────────────────────────────────

\texttt{bootstrap/app.php} is the single configuration file that wires
the application together. It registers middleware aliases, declares the
exception-rendering pipeline, and ensures that all API requests receive
JSON error responses regardless of the exception type.

\begin{lstlisting}[
  style   = php,
  caption = {bootstrap/app.php --- complete file.},
  label   = {lst:bootstrap-full},
]
<?php

use Illuminate\Foundation\Application;
use Illuminate\Foundation\Configuration\Exceptions;
use Illuminate\Foundation\Configuration\Middleware;
use Illuminate\Http\Request;

return Application::configure(basePath: dirname(__DIR__))
    ->withRouting(
        web:      __DIR__.'/../routes/web.php',
        api:      __DIR__.'/../routes/api.php',
        commands: __DIR__.'/../routes/console.php',
        health:   '/up',
    )
    ->withMiddleware(function (Middleware $middleware): void {
        $middleware->alias([
            'clinic.scope' => App\Http\Middleware\ClinicScopeMiddleware::class,
            'role'         => App\Http\Middleware\RoleMiddleware::class,
        ]);
    })
    ->withExceptions(function (Exceptions $exceptions): void {

        // Force JSON responses for all exceptions on api/* routes.
        $exceptions->shouldRenderJsonWhen(
            fn (Request $request) => $request->is('api/*'),
        );

        // 401 --- Unauthenticated (Passport token missing or revoked)
        $exceptions->render(
            function (\Illuminate\Auth\AuthenticationException $_e,
                      Request $_request) {
                return response()->json(
                    ['success' => false, 'message' => 'Token invalid or expired.'],
                    401
                );
            }
        );

        // 422 --- Validation failure (FormRequest)
        $exceptions->render(
            function (\Illuminate\Validation\ValidationException $e,
                      Request $_request) {
                return response()->json(
                    ['success' => false, 'errors' => $e->errors()],
                    422
                );
            }
        );

        // 403 --- Laravel authorization gate failure
        $exceptions->render(
            function (\Illuminate\Auth\Access\AuthorizationException $_e,
                      Request $_request) {
                return response()->json(
                    ['success' => false, 'message' => 'Insufficient permissions.'],
                    403
                );
            }
        );

        // 404 --- Route not found
        $exceptions->render(
            function (
                \Symfony\Component\HttpKernel\Exception\NotFoundHttpException $_e,
                Request $_request
            ) {
                return response()->json(
                    ['success' => false, 'message' => 'Resource not found.'],
                    404
                );
            }
        );

        // 404 --- Model not found (route model binding)
        $exceptions->render(
            function (\Illuminate\Database\Eloquent\ModelNotFoundException $_e,
                      Request $_request) {
                return response()->json(
                    ['success' => false, 'message' => 'Resource not found.'],
                    404
                );
            }
        );

        // Domain exceptions (AppointmentConflictException,
        // InvalidAppointmentStateException, ClinicScopeViolationException,
        // AccountDeactivatedException, SelfDemotionException) are self-
        // rendering via their render() methods and do not require
        // explicit registration here.
    })->create();
\end{lstlisting}

The \texttt{shouldRenderJsonWhen()} callback ensures that even exceptions
not explicitly registered (such as a fatal \texttt{Error} during
development) return a JSON body rather than an HTML stack trace when the
request path begins with \texttt{api/}. This is critical for API clients
that cannot parse HTML error pages.

%% ─────────────────────────────────────────────────────────────────────────
\section{Dependency Injection and the Service Container}
\label{sec:di-container}
%% ─────────────────────────────────────────────────────────────────────────

Laravel's service container resolves every class with type-hinted
constructor parameters automatically. No manual binding is required for
application classes because the container uses PHP reflection to
discover the types of constructor parameters and recursively resolves
their dependencies.

The resolution chain for an incoming request to an assistant appointment
endpoint is:

\begin{enumerate}
  \item The container resolves \phpclass{AssistantAppointmentController}.
    Its constructor requires \phpclass{AppointmentRepository},
    \phpclass{AppointmentService}, and \phpclass{PatientRepository}.
  \item The container resolves \phpclass{AppointmentService}. Its
    constructor requires \phpclass{AppointmentRepository} and
    \phpclass{ScheduleRepository}.
  \item The container resolves all repositories. Their constructors
    have no further dependencies (they rely on static Eloquent model
    methods internally).
  \item The container instantiates the controller with all dependencies
    injected and calls the route-matched action method.
\end{enumerate}

All resolutions happen once per request, not once per method call. The
container caches resolved instances within the request lifecycle, so
both the controller and the service share the same
\phpclass{AppointmentRepository} instance within a single request.

%% ─────────────────────────────────────────────────────────────────────────
\section{Environment Configuration}
\label{sec:environment}
%% ─────────────────────────────────────────────────────────────────────────

The application is configured through \texttt{.env} files following
the twelve-factor app methodology. Every environment-specific value
(database credentials, cache driver, queue driver) is externalised from
the codebase. Table~\ref{tab:env-vars} documents every
environment variable consumed by the application.

\begin{table}[h!]
  \centering
  \caption{Environment variables and their roles.}
  \label{tab:env-vars}
  \renewcommand{\arraystretch}{1.5}
  \begin{tabularx}{\linewidth}{@{}L{4.2cm} L{3cm} X@{}}
    \toprule
    \rowcolor{PrimaryBlue}
    \tblhdr{Variable} & \tblhdr{Example Value} & \tblhdr{Purpose} \\
    \midrule
    \envvar{APP\_KEY}        & base64:...  & Laravel encryption key. Must be 32 random bytes. \\
    \rowalt
    \envvar{APP\_ENV}        & production  & Controls debug output and caching behaviour. \\
    \envvar{DB\_CONNECTION}  & mysql       & Must be \texttt{mysql}; SQLite is not supported. \\
    \rowalt
    \envvar{DB\_HOST}        & 127.0.0.1  & MySQL server hostname. \\
    \envvar{DB\_PORT}        & 3306        & MySQL server port. \\
    \rowalt
    \envvar{DB\_DATABASE}    & cms\_prod   & Database name. \\
    \envvar{DB\_USERNAME}    & cms\_user   & MySQL user with full DML privileges. \\
    \rowalt
    \envvar{DB\_PASSWORD}    & ---         & MySQL password (keep in secret manager). \\
    \envvar{CACHE\_DRIVER}   & redis       & \texttt{redis} for production; \texttt{file} for dev. \\
    \rowalt
    \envvar{QUEUE\_CONNECTION}& redis      & \texttt{redis} for prod; \texttt{sync} for tests. \\
    \envvar{PASSPORT\_...}   & (various)   & OAuth2 client credentials generated by
                                             \texttt{passport:install}. \\
    \bottomrule
  \end{tabularx}
\end{table}

\begin{importantbox}[Test Environment Overrides]
  The \texttt{phpunit.xml} file declares the following overrides for the
  test environment, which take precedence over the \texttt{.env} file
  during PHPUnit test runs:
  \texttt{DB\_CONNECTION=mysql}, \texttt{DB\_DATABASE=CMS},
  \texttt{DB\_USERNAME=root}, \texttt{DB\_PASSWORD=""} (empty string).
  The test suite uses a real MySQL database, not an in-memory SQLite
  substitute. This ensures that MySQL-specific behaviour (FULLTEXT
  indexes, TIME column formatting, UUID storage) is exercised by every
  test run.
\end{importantbox}
